
\documentclass[10pt,landscape,fleqn]{article}
%
\usepackage{pgfplots}
\usepackage[T1]{fontenc}
\usepackage[utf8]{inputenc}
\usepackage[ngerman]{babel}
\usepackage{lmodern}
\usepackage{pgfplots}
\pgfplotsset{compat=1.12}
\usepgfplotslibrary{fillbetween}
\usetikzlibrary{patterns}
%
%\usepackage{pstricks,pst-plot,pst-text,pst-solides3d}
\usepackage{epsfig}
\usepackage{graphicx}
\usepackage[colorlinks=true,
            pdftitle={Mathe I Klausur 03.01.2024},
            pdfauthor={Thomas Carraro},
            pdfpagemode={UseOutlines},pdfstartview={Fit},
            linkcolor=black,urlcolor=blue,
            pdfsubject={Mathematik IV}]{hyperref}
\usepackage{cancel}
%
\usepackage{ifthen}
\newboolean{mitLoes}
%\setboolean{mitLoes}{false}
\setboolean{mitLoes}{true}

\newcounter{sol}
\setcounter{sol}{0}
\ifthenelse{\boolean{mitLoes}}{
\setcounter{sol}{1}
}

\usepackage{amssymb,amsmath,amsthm,am1-aufgaben,bm}
\usepackage{array, hhline}
\usepackage{multirow, multicol}
\usepackage{paralist}
\usepackage{mathtools}

% \usepackage{marvosym,pifont}
% \usepackage{listings}
% \usepackage{algorithmic,algorithm}
% \lstset{language=Matlab,numbers=left}

\newcounter{Aufg}
\setcounter{Aufg}{0}

\newcommand{\Aufgabe}[3][]{
\stepcounter{Aufg}
\subsubsection*{Aufgabe
\arabic{Aufg}}
{#3}
}

\newcommand{\Loesung}[1]{
	\ifthenelse{\boolean{mitLoes}}
	{\subsubsection*{L\"osung:}
		#1}
	{}
	\ifthenelse{\boolean{mitLoes}}{
	\leavevmode
	\newline
	\rule{360pt}{1pt}
\cleardoublepage 
}

}


\textwidth = 270mm
\textheight = 180mm
\leftmargin 5mm
\oddsidemargin = -20mm
\evensidemargin = -20mm
\topmargin = -25mm
\parindent 0cm
\columnsep 2cm

\newcommand{\ds}{\displaystyle}
\newcommand{\rme}{\mathrm{e}}
\newcommand{\rmi}{\mathrm{i}}

\renewcommand{\vec}{\boldsymbol}
\newcommand{\mat}{\boldsymbol}

\newcommand{\ud}{\;\mathrm{d}}
\newcommand{\p}{\partial}
\newcommand{\dO}{\d \vec o}
\newcommand{\dA}{\d \vec A}
\newcommand{\dX}{\d \vec x}

\renewcommand{\div}{\operatorname{div}}
\newcommand{\rot}{\operatorname{rot}}
\newcommand{\grad}{\operatorname{grad}}

\newcommand{\skalar}[1]{\left\langle{#1}\right\rangle}

\renewcommand{\d}{\,{\rm d}}
\newcommand{\x}{\ensuremath{{\vec x}}}
\newcommand{\ve}[3]{\begin{pmatrix} #1\\ #2\\ #3\end{pmatrix}}
\newcommand{\e}{\ensuremath{{\bf e}}}

\begin{document}



\begin{twocolumn}



\href{http://mb-mathe.hsu-hh.de/bause}
{\large\bf Helmut-Schmidt-Universit\"at}\\
\href{http://mb-mathe.hsu-hh.de/bause}
{\large\bf Universit\"at der Bundeswehr Hamburg}\\
\href{http://mb-mathe.hsu-hh.de/bause}
{\large\bf Fakult\"at f\"ur Maschinenbau und \\ Bauingenieurwesen}\\[1ex]
\href{http://mb-mathe.hsu-hh.de/bause}
{\large\bf Prof. Dr.\ Thomas Carraro}\\
\href{http://mb-mathe.hsu-hh.de/bause}
{\large\bf Dr.\ Ulrike Kochan-Eilers}\\
\\[-28mm]
\phantom{.}
\hfill\includegraphics[height=20mm]{HSU_RGB}

\vspace*{8mm}
{\Large\bf Mathematik III/B (WI/ET) \hfill Zusatzblatt}\\[1ex]
{\Large FT 2025 \hfill  Integrationsgrenzen}


\vspace*{1mm}
\rule{360pt}{2pt}

\vspace*{0ex}
% \paragraph*{Ex1}
%
\renewcommand{\d}{\mathrm{d}}
\renewcommand{\vec}[1]{\mathbf{#1}}
\newcommand{\EH}[1]{\mathrm{e}^{#1}}

\vspace*{+0.25cm}
Geben Sie die Integrationsgrenzen für die folgenden Körper $K$ in geeigneten Koordinaten an.

\begin{abc}
\item $K = \left\{\textbf{x}\in  \mathbb{R}^2: 4 \leq x^2+y^2 \leq 9 \right\} $
\item   $ K = \left\{\textbf{x}\in  \mathbb{R}^2: 4 \leq x^2+y^2 \leq 9, \,  y \leq 0 \right\}$

\item $ K = \left\{\textbf{x}\in  \mathbb{R}^3:  x^2+y^2 \leq 9, \, z \geq 0 \right\} $
\item $ K = \left\{\textbf{x}\in  \mathbb{R}^3:  x^2+y^2 \leq 9, \, y\geq 0, \, z\geq 0 \right\} $
\item $ K = \left\{\textbf{x}\in  \mathbb{R}^3:  x^2+y^2 \leq 9, \, 0\leq z\leq24-2x-8y \right\} $
\item $ K = \left\{\textbf{x}\in  \mathbb{R}^3:  x^2+y^2 \leq 9, \, x \leq 0, \, y\leq0, z\geq 0 \right\} $
\item $ K = \left\{\textbf{x}\in  \mathbb{R}^3:  x^2+y^2+z^2 \leq 9\right\} $
\item $K = \left\{\textbf{x}\in  \mathbb{R}^3:  4\leq x^2+y^2+z^2 \leq 9, \, z\geq0\right\} $
\item $ K = \left\{\textbf{x}\in  \mathbb{R}^3:  4\leq x^2+y^2+z^2 \leq 9, \, z\geq0, y\geq0 \right\} $
\item $ K = \left\{\textbf{x}\in  \mathbb{R}^3:  x^2+y^2+z^2 \leq 9, \, x\leq0, \, y\geq0 \right\} $
\item $ K = \left\{\textbf{x}\in  \mathbb{R}^3:  x^2+y^2+z^2 \leq 4, \, x\geq 0, \, y\leq0, z\leq0 \right\} $
\end{abc}

\vspace*{+0.25cm}
\newpage


\ifthenelse{\boolean{mitLoes}}{
\textbf{L\"osung:}\\
\begin{abc}
\item $$K=\left\{\textbf{x} \in \mathbb{R}^2 :\,0 \leq \varphi \leq 2\pi,\, 
2\leq r \leq 3 \right\}$$
\item $$K=\left\{\textbf{x} \in \mathbb{R}^2 :\, \pi \leq \varphi \leq 2\pi,\, 
2\leq r \leq 3 \right\}$$
\item $$K=\left\{\textbf{x} \in \mathbb{R}^3 :\, 0 \leq \varphi \leq 2\pi,\, 
0\leq r \leq 3, 0\leq z\leq \infty \right\}$$
\item $$K=\left\{\textbf{x} \in \mathbb{R}^3 :\, 0 \leq \varphi \leq \pi,\, 
0\leq r \leq 3, 0\leq z\leq \infty \right\}$$
\item $$K=\left\{\textbf{x} \in \mathbb{R}^3 :\, 0 \leq \varphi \leq 2\pi,\, 
0\leq r \leq 3, 0\leq z\leq 24-2r\cos \varphi -8r \sin \varphi\right\}$$
\item $$K=\left\{\textbf{x} \in \mathbb{R}^3 :\, \pi \leq \varphi \leq \frac{3}{2}\pi,\, 0\leq r \leq 3, 0\leq z\leq \infty \right\}$$
\item $$K=\left\{\textbf{x} \in \mathbb{R}^3 :\, 0 \leq \varphi \leq 2\pi,\, 
0\leq r \leq 3, 0 \leq \theta \leq \pi \right\}$$
\item $$K=\left\{\textbf{x} \in \mathbb{R}^3 :\, 0 \leq \varphi \leq 2\pi,\, 
2\leq r \leq 3, 0\leq \theta \leq \frac{\pi}{2} \right\}$$
\item $$K=\left\{\textbf{x} \in \mathbb{R}^3 :\, 0 \leq \varphi \leq \pi,\, 
2\leq r \leq 3, 0\leq \theta \leq \frac{\pi}{2} \right\}$$
\item $$K=\left\{\textbf{x} \in \mathbb{R}^3 :\, \frac{\pi}{2} \leq \varphi \leq \pi,\, 0\leq r \leq 3, 0\leq \theta \leq \pi \right\}$$
\item $$K=\left\{\textbf{x} \in \mathbb{R}^3 :\, \frac{3}{2} \pi \leq \varphi \leq 2\pi,\, 0\leq r \leq 2, \frac{\pi}{2} \leq \theta \leq \pi \right\}$$
\end{abc}
}{}

% 



%%%%%%%%%%%%%%%%%%%%%%%%%%%%%%%%%%%%%%%%%%%%%%%%%%%%%%%%%%%%%%%%%%%%%%%%%%%%%%%%


%%%%%%%%%%%%%%%%%%%%%%%%%%%%%%%%%%%%%%%%%%%%%%%%%%%%%%%%%%%%%%%%%%%%%%%%%%%%%%%%
\end{twocolumn}

\end{document}

%%%%%%%%%%%%%%%%%%%%%%%%%%%%%%%%%%%%%%%%%%%%%%%%%%%%%%%%%%%%%%%%%%%%%%%%%%%%%%%%
